\chapter[1987 --- Tonegawa]{1987: Susumu Tonegawa}
\label{ch:1987_susumutonegawa}

\section*{Main Contribution}

\textit{Discovered the genetic mechanism (V(D)J recombination) that generates the enormous diversity of antibodies, explaining how the immune system can recognize virtually any pathogen.}

\section{Official Citation}

for his discovery of the genetic principle for generation of antibody diversity

The 1987 Nobel Prize in Physiology or Medicine was awarded to Susumu Tonegawa ``for his discovery of the genetic principle for generation of antibody diversity.'' Tonegawa's work solved one of the most perplexing puzzles in immunology: how the human genome, with its limited number of genes, can produce the vast diversity of antibody molecules needed to recognize millions of different foreign antigens. His discovery of somatic recombination---the process by which antibody genes are assembled from smaller gene segments during the development of B lymphocytes---revealed a fundamentally new mechanism of gene expression and transformed our understanding of how the immune system generates its notable repertoire of antigen-binding specificities.

\section{Background and Historical Context}

The immune system's ability to recognize and respond to millions of different foreign molecules (antigens) is one of a notable features of vertebrate biology. Antibodies---the Y-shaped proteins produced by B lymphocytes that mediate humoral immunity---are the molecular tools through which this recognition is achieved. Each antibody molecule has two identical antigen-binding sites (paratopes) that recognize a specific molecular feature (epitope) on the surface of the antigen. The human immune system can produce antibodies that recognize virtually any molecular structure, from the simplest chemical group to the most complex protein or polysaccharide.

By the 1960s, it was known that antibodies are composed of two types of polypeptide chains: heavy chains and light chains. Each antibody molecule consists of two identical heavy chains and two identical light chains, linked by disulfide bonds. The antigen-binding site is formed by the variable (V) regions of the heavy and light chains, while the constant (C) regions determine the antibody's effector function (its ability to activate complement, bind to Fc receptors, and so on).

The central puzzle of immunology was the ``diversity problem.'' Estimates of the number of different antibody specificities that the human immune system can produce ranged from $10^6$ to $10^9$ or more. However, the entire human genome contains only about 20,000--30,000 genes, and a single gene encodes each polypeptide chain. How could a limited number of genes give rise to such an enormous diversity of antibody molecules?

Two principal hypotheses were proposed to explain antibody diversity. The ``germ-line hypothesis'' held that the diversity of antibodies is encoded directly in the germ-line DNA---that is, that each antibody specificity is encoded by a separate gene, and that the genome contains a very large number of antibody genes. The ``somatic mutation hypothesis'' held that antibody diversity is generated by somatic (non-inheritable) mutations that occur during the development of B lymphocytes, creating new antibody specificities that were not present in the germ-line DNA.

Susumu Tonegawa was born in Nagoya, Japan, in 1939, and trained in chemistry at Kyoto University. He emigrated to the United States in 1963 to pursue graduate studies at the University of California, San Diego, where he earned his Ph.D. in molecular biology in 1968. He then joined the laboratory of Richard Dulbecco at the Salk Institute in La Jolla, California, where he began studying the genetic basis of antibody diversity.

\section{The Discovery/Contribution}

Tonegawa's approach to the diversity problem was ingenious. Rather than trying to measure the diversity of antibodies directly (a technically daunting task), he decided to compare the organization of antibody genes in DNA from antibody-producing cells (B lymphocytes) with the organization of the same genes in DNA from non-antibody-producing cells (such as embryonic cells or sperm). If the diversity of antibodies was encoded directly in the germ-line DNA (as the germ-line hypothesis predicted), then the organization of antibody genes should be the same in both types of cells. If, on the other hand, the antibody genes were rearranged during B lymphocyte development (as the somatic mutation hypothesis suggested), then the organization of antibody genes should be different in antibody-producing cells compared to non-antibody-producing cells.

Tonegawa designed an experiment to test this prediction. Working with his colleague Nobuhiko Hozumi in the early 1970s, Tonegawa used restriction enzymes and DNA hybridization to compare the organization of antibody genes in DNA from two types of cells: myeloma cells (cancerous B lymphocytes that produce large quantities of a single antibody) and embryonic cells (which do not produce antibodies). They used radioactive probes complementary to the sequences encoding the variable (V) and constant (C) regions of antibody heavy chains, and they hybridized these probes to restriction fragments of DNA from the two cell types.

The results were striking. Tonegawa and Hozumi found that the V and C gene segments were far apart in the DNA of embryonic cells---they were located on different restriction fragments, separated by thousands of base pairs. However, in myeloma cells, the V and C gene segments were close together---they were located on the same restriction fragment. This demonstrated that the antibody genes had been rearranged during B lymphocyte development: the V and C gene segments, which were separated in the germ-line DNA, had been brought together by a DNA rearrangement event to form a complete, functional antibody gene.

Tonegawa published his landmark paper in 1976, and it immediately attracted worldwide attention. The discovery of DNA rearrangement in antibody genes---later called V(D)J recombination---was a paradigm-shifting event in biology. It demonstrated that the genome is not a static, fixed structure, but rather a dynamic system in which gene segments can be rearranged to produce new genetic information.

\subsection{The Mechanism of V(D)J Recombination}

Subsequent research, by Tonegawa and many others, elucidated the molecular mechanism of V(D)J recombination. The antibody heavy chain gene is assembled from three types of gene segments: variable (V) segments, diversity (D) segments, and joining (J) segments. In the germ-line DNA, these segments are present in multiple copies---for example, the human heavy chain locus contains approximately 40 functional V$_H$ segments, 25 D segments, and 6 J segments. During B lymphocyte development, one V, one D, and one J segment are randomly selected and joined together by a recombination process that involves the RAG1 and RAG2 (recombination-activating gene) enzymes. The resulting VDJ exon encodes the variable region of the heavy chain.

The light chain gene is assembled by a similar process, using V and J segments (but no D segments). The combination of different V, D, and J segments in the heavy and light chain genes generates a large number of possible antibody specificities through a mechanism known as combinatorial diversity.

In addition to combinatorial diversity, several other mechanisms contribute to antibody diversity: junctional diversity (the addition or deletion of nucleotides at the junctions between gene segments during recombination), somatic hypermutation (the introduction of point mutations into the variable region genes of antibody molecules during the immune response), and combinatorial association of different heavy and light chains. Together, these mechanisms generate an estimated $10^{11}$ or more different antibody specificities---far more than could be encoded directly in the germ-line DNA.

\section{Impact on Modern Medicine}

Tonegawa's discovery of V(D)J recombination has had a significant impact on modern medicine and biology. The most immediate practical consequence was the elucidation of the genetic basis of certain immune deficiencies. Mutations in the RAG1 or RAG2 genes, which are essential for V(D)J recombination, cause severe combined immunodeficiency (SCID)---a life-threatening condition in which the patient lacks functional B and T lymphocytes. Understanding the genetic basis of SCID has enabled the development of genetic tests for the disease and has paved the way for gene therapy approaches to its treatment.

The study of V(D)J recombination has also provided insights into the mechanisms of lymphoid malignancy. Chromosomal translocations involving the antibody gene loci---in which segments of the antibody genes are joined to proto-oncogenes on other chromosomes---are a major cause of B cell lymphomas and leukemias. For example, the translocation t(8;14), which joins the MYC oncogene to the immunoglobulin heavy chain locus, is a hallmark of Burkitt lymphoma. Understanding these translocations has provided insights into the mechanisms of lymphoid malignancy and has suggested new approaches to treatment.

Tonegawa's discovery has also been essential for the development of therapeutic antibodies and immunotherapies. The understanding of antibody gene organization and V(D)J recombination has enabled the engineering of antibodies with desired specificities and properties, including humanized antibodies for cancer therapy and bispecific antibodies for the treatment of infectious diseases and cancer.

In basic biology, the discovery of V(D)J recombination has revealed a fundamentally new mechanism of gene expression: the assembly of functional genes from smaller gene segments by regulated DNA rearrangement. This mechanism is not unique to antibody genes; similar processes of DNA rearrangement occur in the generation of T cell receptors, and the general principle of gene segment rearrangement has been found in other contexts as well. The study of V(D)J recombination has thus contributed to our understanding of genome organization, gene regulation, and the evolution of immune systems.

Tonegawa's work has also influenced the field of synthetic biology. The recognition that the immune system generates diversity through the regulated rearrangement of gene segments has inspired the development of synthetic systems for generating protein diversity, including phage display libraries and other in vitro evolution technologies. These technologies are now widely used in drug discovery, diagnostics, and basic research.

\section{Legacy}

Susumu Tonegawa continued his work on the molecular biology of the immune system at the Massachusetts Institute of Technology (MIT), where he founded the Center for Neurobiology and Memory. In a surprising pivot, Tonegawa later turned his attention to the neurobiology of memory, studying the molecular mechanisms by which memories are formed and stored in the brain. He has made important contributions to our understanding of the role of the hippocampus in memory formation and the mechanisms of long-term potentiation (LTP), a cellular model of learning and memory.

Tonegawa's career illustrates a remarkable breadth of scientific curiosity and achievement. From his Nobel Prize-winning work on antibody diversity to his current research on the neurobiology of memory, Tonegawa has consistently pushed the boundaries of biological knowledge and has demonstrated that a major discoveries often come from asking fundamental questions about how nature works.

The discovery of V(D)J recombination remains one of the major findings in the history of biology. It solved one of the central puzzles of immunology, revealed a new mechanism of gene expression, and provided tools and concepts that continue to influence medicine, biology, and biotechnology. Susumu Tonegawa's work is a testament to the power of elegant experimental design, rigorous scientific thinking, and the willingness to challenge established dogma.


\section{Modern Developments}

Tonegawa's antibody diversity work has led to modern immunogenomics. Single-cell sequencing reveals the repertoire of B and T cells in health and disease. Understanding of V(D)J recombination has informed development ofCAR-T cell therapy. Understanding of somatic hypermutation has enabled development of broadly neutralizing antibodies for HIV and influenza. Computational prediction of antibody-antigen interactions accelerates therapeutic antibody development.


\section{Bibliography}

\begin{thebibliography}{9}
\bibitem{nobel-1987}
Nobel Prize Outreach, ``The Nobel Prize in Physiology or Medicine 1987,''\emph{NobelPrize.org}.\\
\url{https://www.nobelprize.org/prizes/medicine/1987/summary/}

\bibitem{lecture-1987}
Susumu Tonegawa, ``Nobel Lecture,''\emph{NobelPrize.org}. Winner-authored primary source; downloadable from the official Nobel Prize archive.\\
\url{https://www.nobelprize.org/prizes/medicine/1987/tonegawa/lecture/}
\end{thebibliography}
